\DocumentMetadata{
  pdfversion=2.0,
  pdfstandard=ua-2,
  lang=en-US,
  tagging=on,
  tagging-setup={math/setup=mathml-SE,tabsorder=structure}
}
\documentclass[11pt,letterpaper]{article}
\usepackage[margin=0.68in,includefoot]{geometry}
\usepackage{newcomputermodern}
\usepackage{amsmath,amscd,mathtools}
\usepackage{tikz,adjustbox}
\providecommand{\square}{\mdlgwhtsquare}
\usepackage[hidelinks]{hyperref}
\hypersetup{
  pdftitle={Complex Analysis PhD Exam, September 1992},
  pdfauthor={University of Florida Department of Mathematics},
  pdfsubject={Graduate mathematics examination},
  pdfdisplaydoctitle=true,
  bookmarksopen=true
}
\setcounter{secnumdepth}{0}
\setlength{\parindent}{0pt}
\setlength{\parskip}{0.28em}
\setlength{\emergencystretch}{3em}
\setlength{\leftmargini}{2.15em}
\setlength{\leftmarginii}{2.65em}
\setlength{\leftmarginiii}{3em}
\pagestyle{empty}
\raggedbottom
\allowdisplaybreaks
\NewTaggingSocketPlug{block/list/label}{examproblem}{%
  \tagstructbegin{tag=Lbl}\tagstructend%
  \tagstructbegin{tag=\LBody}%
  \tagstructbegin{tag=\ProblemHeadingTag,title-o={\ProblemHeadingText}}%
  \tagmcbegin{tag=\ProblemHeadingTag,actualtext={\ProblemHeadingText}}%
  #2%
  \tagmcend%
  \tagstructend%
}
\newenvironment{examproblems}[2]{%
  \def\ProblemHeadingTag{#2}%
  \AssignTaggingSocketPlug{block/list/label}{examproblem}%
  \begin{enumerate}%
  \setcounter{enumi}{#1}%
  \renewcommand{\labelenumi}{\textbf{\theenumi.}}%
  \setlength{\topsep}{0.35em}%
  \setlength{\partopsep}{0pt}%
  \setlength{\itemsep}{0.48em}%
  \setlength{\parsep}{0.18em}%
}{\end{enumerate}}
\newenvironment{examsubparts}{%
  \AssignTaggingSocketPlug{block/list/label}{default}%
  \begin{enumerate}%
  \setlength{\topsep}{0.18em}%
  \setlength{\partopsep}{0pt}%
  \setlength{\itemsep}{0.16em}%
  \setlength{\parsep}{0.1em}%
}{\end{enumerate}}
\newenvironment{examinstructions}{%
  \par\begingroup\itshape
}{%
  \par\endgroup\vspace{0.18em}
}
\makeatletter
\renewcommand\subsection{\@startsection{subsection}{2}{0pt}%
  {-1.25ex plus -.25ex minus -.1ex}%
  {0.55ex plus .1ex}%
  {\normalfont\large\bfseries}}
\renewcommand{\@maketitle}{%
  \newpage\null\vskip -1em%
  {\footnotesize\bfseries
    \href{https://www.ufl.edu/}{University of Florida}, \href{https://math.ufl.edu/}{Department of Mathematics}\par}%
  \vskip -0.5em%
  \tagstructbegin{tag=H1,title={Complex Analysis PhD Exam, September 1992}}%
  \tagpdfparaOff%
  \tagmcbegin{tag=H1}%
    {\LARGE\bfseries\href{https://gma.math.ufl.edu/exams/complex-analysis-phd/}{Complex Analysis PhD Exam}, September 1992\par}%
  \tagmcend%
  \tagpdfparaOn%
  \tagstructend%
  \vskip 0.25em%
  \tagmcbegin{artifact}%
    \hrule height 1pt%
  \tagmcend%
  \vskip 1em%
}
\makeatother
\title{Complex Analysis PhD Exam, September 1992}
\author{}
\date{}
\begin{document}
\maketitle
\thispagestyle{empty}
\begin{examinstructions}
Present the solutions with the necessary details, so the partial credit can be properly assessed. Do all 9 problems.
\end{examinstructions}
\begin{examproblems}{0}{H2}
\def\ProblemHeadingText{Problem 1.}
\item
Construct a conformal map which maps the vertical strip \(0<\operatorname{Re} z<1\) onto the upper half-plane.
\def\ProblemHeadingText{Problem 2.}
\item
The function
\[
f(z)=\sum_{n=1}^{\infty}\frac{z^n}{n^2}
\]
 is analytic in the unit circle \(|z|<1\). Can it be continued analytically beyond the circle \(|z|=1\)? Discuss the nature of the singularities (if any) on \(|z|=1\).
\def\ProblemHeadingText{Problem 3.}
\item
This problem has two parts.
\begin{examsubparts}
\item[\textnormal{(a)}]
Show that the function
\[
f(z)=\sum_{n=1}^{\infty}\frac{(-1)^{n+1}}{n^z}
\]
 is analytic in \(\operatorname{Re} z>0\).
\item[\textnormal{(b)}]
Show that it can be analytically continued to the entire complex plane.
\end{examsubparts}
\def\ProblemHeadingText{Problem 4.}
\item
Let \(f(z)\) be analytic in the disc \(|z|\leq 1\). If
\[
|f(z)|\leq 1 \quad\text{for }|z|\leq 1,
\]
 prove that
\[
\frac{|f'(z)|}{1-|f(z)|^2}\leq\frac{1}{1-|z|^2}\quad\text{for }|z|\leq 1.
\]
\def\ProblemHeadingText{Problem 5.}
\item
Let \(F(z)\) be an entire function with the property that \(\operatorname{Re}F(z)\) is bounded from above. Prove that \(F(z)\) is constant.
\def\ProblemHeadingText{Problem 6.}
\item
Let~\(\Omega\) be an open connected subset of~\(\mathbb{C}\). Let \(\{f_n\}\) be a sequence of one-to-one analytic functions on~\(\Omega\), and assume that \(\{f_n\}\) converges uniformly on any compact subset of~\(\Omega\) to a function~\(f\). Show that~\(f\) is either one-to-one or a constant, and show that both possibilities can occur.
\def\ProblemHeadingText{Problem 7.}
\item
Suppose that~\(f\) is entire, of order~\(\rho\). Let \(n(r)\) be the number of zeros (counting the multiplicity) of~\(f\) in the disk \(\{z:|z|\leq r\}\). Prove that
\[
\limsup_{r\to\infty}\frac{\log n(r)}{\log r}\leq\rho,
\]
 and show, by examples, that equality can happen.
\def\ProblemHeadingText{Problem 8.}
\item
The expression
\[
\{f,z\}=\frac{f'''(z)}{f'(z)}-\frac{3}{2}\left\{\frac{f''(z)}{f'(z)}\right\}^2
\]
 is called the Schwarzian derivative of~\(f\). Prove the following:
\begin{examsubparts}
\item[\textnormal{(a)}]
If \(f(z)=a(z-z_0)^m+\cdots\), where~\(m\) is an integer, then
\[
\{f,z\}=A(z-z_0)^{-2}+\cdots.
\]
\item[\textnormal{(b)}]
Find the value of~\(A\).
\item[\textnormal{(c)}]
What can you say about \(\{f,z\}\) at the point \(z_0\) if \(m=\pm1\)?
\end{examsubparts}
\def\ProblemHeadingText{Problem 9.}
\item
This problem has two parts.
\begin{examsubparts}
\item[\textnormal{(a)}]
Evaluate
\[
\frac{1}{2\pi i}\int_{2-i\infty}^{2+i\infty}\frac{x^z}{z(z+1)}\,dz
\]
 for \(x>1\).

\par
Hint: Write the above integral as \(\lim_{t\to\infty}I(t)\), where
\[
I(t)=\frac{1}{2\pi i}\int_{2-it}^{2+it}\frac{x^z}{z(z+1)}\,dz,
\]
 and evaluate \(I(t)\) by using the residue theorem on the rectangle \(2-it\), \(2+it\), \(b+it\), and \(b-it\), where~\(b\) is a large negative number.
\item[\textnormal{(b)}]
What happens if \(0<x<1\)?
\end{examsubparts}
\end{examproblems}
\end{document}
